

\SGARefTarget{sga3:iv:subsection:4:7}\subsection*{4.7. Use of Effectivity Criteria: Isomorphism Theorem}
\addcontentsline{toc}{subsection}{\SGARefLink{sga3:iv:subsection:4:7}{4.7}. Use of Effectivity Criteria: Isomorphism Theorem}

In this subsection we give an example of the use of effectivity criteria.
The initial data are a topology \(\mathcal T\) on \(\mathscr C\), still
assumed coarser than the canonical topology; a family \((M)\) of morphisms
of \(\mathscr C\) satisfying Axioms \SGARefLink{sga3:iv:para:4:6:3:item:a:01}{(a)}--\SGARefLink{sga3:iv:para:4:6:3:item:e:mathcal:t:01}{\((e_{\mathcal T})\)} of \SGARefLink{sga3:iv:para:4:6:3}{4.6.3};
and a family \((N)\) of morphisms of \(\mathscr C\) which may satisfy the
following axioms:
\begin{enumerate}
  \item[(a)]\SGARefTarget{sga3:iv:subsection:4:7:item:a:01} \((N)\) is stable under base change;
  \item[\((f_{\mathcal T})\)]\SGARefTarget{sga3:iv:subsection:4:7:item:f:mathcal:t:01} morphisms in \((N)\) descend for the given
    topology. More precisely, for every
    \(S\in\operatorname{Ob}(\mathscr C)\), every covering family
    \(\{S_i\to S\}\), and every family \(\{T_i\to S_i\}\) of morphisms
    in \((N)\), every descent datum on the \(T_i\) relative to
    \(\{S_i\to S\}\) is effective; if \(T\) denotes the descended
    object, then \(T\to S\) belongs to \((N)\).
\end{enumerate}

Since every element of \((M)\) is covering (condition \SGARefLink{sga3:iv:para:4:6:3}{4.6.3}
\SGARefLink{sga3:iv:para:4:6:3:item:d:mathcal:t:01}{\((d_{\mathcal T})\)}), Axiom \SGARefLink{sga3:iv:subsection:4:7:item:f:mathcal:t:01}{\((f_{\mathcal T})\)} implies \SGARefLink{sga3:iv:subsection:4:7:item:f:m:01}{the following
axiom}:\footnote{Editor's note: Axiom \SGARefLink{sga3:iv:subsection:4:7:item:f:m:01}{\((f_M)\)}, which originally appeared
before Proposition \SGARefLink{sga3:iv:proposition:4:7:2}{4.7.2}, has been moved here.}
\begin{enumerate}
  \item[\((f_M)\)]\SGARefTarget{sga3:iv:subsection:4:7:item:f:m:01} if \(Y\to X\) belongs to \((N)\) and
    \(X\to X_1\) belongs to \((M)\), every descent datum on \(Y\)
    relative to \(X\to X_1\) is effective; if \(Y_1\) denotes the
    descended object, then \(Y_1\to X_1\) belongs to \((N)\).
\end{enumerate}

Here is an example, to be treated later: \(\mathscr C\) is the category of
schemes, \(\mathcal T\) is the faithfully flat quasi-compact topology,
\((M)\) is the family of faithfully flat quasi-compact morphisms, and
\((N)\) is the family of closed immersions, or the family of quasi-compact
immersions.

Taking \SGARefLink{sga3:iv:remark:4:6:11}{4.6.11} into account, the principal result of \SGARefLink{sga3:iv:theorem:4:6:6}{4.6.6} may be recalled
as follows.

\medskip
\noindent\SGARefTarget{sga3:iv:proposition:4:7:1}\textbf{Proposition 4.7.1.}
\textit{Let \(X\) be a squarable object of \(\mathscr C\), and let \(R\)
be an equivalence relation of type \((M)\) in \(X\). Suppose that
\(R\to X\times X\) belongs to \((N)\), where \((N)\) satisfies \SGARefLink{sga3:iv:subsection:4:7:item:a:01}{(a)} and
\SGARefLink{sga3:iv:subsection:4:7:item:f:mathcal:t:01}{\((f_{\mathcal T})\)}. Then the sheaf quotient \(X/R\) is given by}
\[
\begin{aligned}
 (X/R)(S)=\{Z\subseteq X_S\mid{}&
   Z\text{ is an }S\text{-subobject stable under }R\times S,\\
 & Z\to X_S\text{ belongs to }(N),\quad
   Z\to S\text{ is covering or belongs to }(M),\\
 & R_Z\simeq Z\times_S Z\,\}.
\end{aligned}
\]

We also have \SGARefLink{sga3:iv:proposition:4:7:2}{the following result}.

\medskip
\noindent\SGARefTarget{sga3:iv:proposition:4:7:2}\textbf{Proposition 4.7.2.}
\textit{Let \(X\in\operatorname{Ob}(\mathscr C)\), let \(R\) be an
\((M)\)-effective equivalence relation in \(X\), and let \((N)\) be a
family of morphisms satisfying \SGARefLink{sga3:iv:subsection:4:7:item:a:01}{(a)} and \SGARefLink{sga3:iv:subsection:4:7:item:f:m:01}{\((f_M)\)}. For every subobject
\(Y\) of \(X\) which is stable under \(R\) and for which \(Y\to X\)
belongs to \((N)\), the equivalence relation \(R_Y\) induced in \(Y\) by
\(R\) is \((M)\)-effective. Moreover, the quotient}
\[
  Y/R_Y=Y'
\]
\textit{is a subobject of \(X'=X/R\), and \(Y'\to X'\) belongs to
\((N)\).}

\textit{The map \(Y\mapsto Y'\) is a bijection from the set of subobjects
\(Y\) of \(X\) which are stable under \(R\) and for which \(Y\to X\)
belongs to \((N)\), onto the set of subobjects \(Y'\) of \(X'\) for which
\(Y'\to X'\) belongs to \((N)\). Its inverse is}
\[
  Y'\longmapsto Y'\times_{X'}X.
\]

Indeed, since \(R\) is \((M)\)-effective, the morphism \(X\to X'\)
belongs to \((M)\). Let \(Y'\) be a subobject of \(X'\) such that the
canonical morphism \(Y'\to X'\) belongs to \((N)\). Then
\(Y=Y'\times_{X'}X\) is a subobject of \(X\) stable under \(R\), and the
morphisms \(Y\to X\) and \(Y\to Y'\) belong respectively to \((N)\) and
\((M)\), since both families are stable under base change. Let \(R_Y\)
denote the equivalence relation induced in \(Y\) by \(R\). By \SGARefLink{sga3:iv:proposition:4:4:11}{4.4.11}, the
sheaf quotient \(Y/R_Y\) is represented by \(Y'\); hence \SGARefLink{sga3:iv:lemma:4:6:5}{4.6.5} shows that
\(R_Y\) is \((M)\)-effective.

Conversely, let \(Y\) be a subobject of \(X\), stable under \(R\), such
that \(Y\to X\) belongs to \((N)\). Since \(Y\) is stable under \(R\),
its two inverse images in
\[
  R=X\times_{X'}X
\]
are identical, and \(Y\) is equipped with a descent datum relative to
\(X\to X'\). The required conclusion follows from Axiom \SGARefLink{sga3:iv:subsection:4:7:item:f:m:01}{\((f_M)\)} for
the family \((N)\).

\medskip
\noindent\SGARefTarget{sga3:iv:corollary:4:7:3}\textbf{Corollary 4.7.3.}
\textit{Let \(X\in\operatorname{Ob}(\mathscr C)\), and let \(R\) be an
\((M)\)-effective equivalence relation in \(X\). Suppose in addition that
\(R\to X\times X\) belongs to \((N)\), where \((N)\) satisfies \SGARefLink{sga3:iv:subsection:4:7:item:a:01}{(a)} and
\SGARefLink{sga3:iv:subsection:4:7:item:f:mathcal:t:01}{\((f_{\mathcal T})\)}. Then, for every \(Y\) as in \SGARefLink{sga3:iv:proposition:4:7:2}{4.7.2},
\(R_Y\to Y\times Y\) also belongs to \((N)\), and consequently \SGARefLink{sga3:iv:proposition:4:7:1}{4.7.1}
gives}
\[
\begin{aligned}
 (Y/R_Y)(S)=\{Z\subseteq Y_S\mid{}&
   Z\text{ is an }S\text{-subobject stable under }R_Y\times S,\\
 & Z\to Y_S\text{ belongs to }(N)
   \text{ (hence so does }Z\to X_S),\\
 & Z\to S\text{ is covering},\quad R_Z\simeq Z\times_S Z\,\}.
\end{aligned}
\]
